% =====================================================================
\parte{Proyecto de fase}{cap:proyecto}
% =====================================================================

\section*{Perplejidad: modelo uniforme, bigramas y trigramas}

\paragraph{Objetivo.} Medir, con la métrica con la que se comparan los modelos de lenguaje, cuánto mejora un modelo de caracteres al ampliar el contexto, y explicar por qué más contexto no siempre es mejor cuando hay pocos datos. El proyecto integra el modelo de $n$-gramas y el suavizado (M01, §2.3 y~§3.3), la perplejidad (M01, §5.1), los tensores (M03, §6.1) y la conversión a bits (M06, §2.3).

\paragraph{Datos.} La lista de 316 nombres del cuaderno de prácticas del M01, con el mismo vocabulario de 28 tokens (\fin\ y las 27 letras del alfabeto español) y la misma partición: generador \texttt{np.random.default\_rng(7)}, un $15\,\%$ para validación (47 nombres, 321 tokens que predecir) y los 269 restantes para entrenamiento. No se descarga ningún dato.

\paragraph{Plazo.} Una semana desde la realización del examen. El proyecto cuenta en la nota de la fase, no en la del examen.

\begin{figure}[h]
\centering
\begin{tikzpicture}[node distance=5mm and 7mm,
  caja/.style={draw,rounded corners=2pt,align=center,font=\small,minimum height=8mm,inner sep=3pt},
  flecha/.style={-{Stealth[length=2mm]}}]
\node[caja] (n) {Nombres};
\node[caja,right=of n] (p) {Partición\\entrenamiento / validación};
\node[caja,right=of p] (b) {Bigramas\\$N\in\R^{28\times 28}$};
\node[caja,above=of b] (u) {Uniforme};
\node[caja,below=of b] (t) {Trigramas\\$N_3\in\R^{28\times28\times28}$};
\node[caja,right=of b] (s) {Suavizado $\alpha$\\interpolación $\lambda$};
\node[caja,right=of s] (g) {Pérdida y\\perplejidad\\de validación};
\draw[flecha] (n)--(p);
\draw[flecha] (p)--(b); \draw[flecha] (p.north) |- (u); \draw[flecha] (p.south) |- (t);
\draw[flecha] (b)--(s); \draw[flecha] (t) -| (s);
\draw[flecha] (s)--(g); \draw[flecha] (u) -| (g);
\end{tikzpicture}
\par\smallskip{\small Flujo de trabajo del proyecto de fase.}
\end{figure}

\paragraph{Entregables.}
\begin{enumerate}[label=\arabic*.]
  \item Un cuaderno \texttt{H1-proyecto.ipynb} que se ejecute de principio a fin en CPU en menos de un minuto y sin errores.
  \item Una tabla con la pérdida de validación (en nats y en bits) y la perplejidad de cada modelo.
  \item Una gráfica de la perplejidad de validación frente a $\alpha$, con escala logarítmica en $\alpha$, para bigramas y trigramas.
  \item Un informe breve (a lo sumo una página, en celdas de texto del cuaderno) que responda a las preguntas del paso~6.
\end{enumerate}

\paragraph{Pasos.}
\begin{enumerate}[label=\arabic*.]
  \item \emph{Tokenización con relleno.} Para bigramas, cada nombre se rodea de un símbolo \fin\ a cada lado; para trigramas, se antepone \fin\fin\ y se añade \fin\ al final. Así, ambos modelos predicen exactamente los mismos $\ell+1$ tokens de un nombre de $\ell$ letras. Debe comprobarse con un \texttt{assert} que el número de predicciones coincide.
  \item \emph{Modelo uniforme.} $P(x_t)=1/28$. Justifique, sin hacer ningún cálculo sobre los datos, que su perplejidad es $28$.
  \item \emph{Bigramas.} Con la matriz $N$ y el suavizado de Laplace generalizado del M01, $P_\alpha(j\mid i)=(N_{ij}+\alpha)/\big(\sum_k N_{ik}+28\alpha\big)$.
  \item \emph{Trigramas.} Con el tensor $N_3$ de $28\times 28\times 28$ conteos, $P_\alpha(k\mid i,j)=(N_{3,ijk}+\alpha)/\big(\sum_l N_{3,ijl}+28\alpha\big)$. Cuente cuántas casillas son distintas de cero.
  \item \emph{Barrido de $\alpha$.} Pruebe al menos $\alpha\in\{0{,}01;\ 0{,}05;\ 0{,}1;\ 0{,}2;\ 0{,}5;\ 1\}$ en ambos modelos y elija $\alpha$ en validación. Se recomienda separar una tercera partición de prueba para que la elección de $\alpha$ no contamine la medida final.
  \item \emph{Análisis.} (a) ¿Qué modelo obtiene la menor perplejidad y con qué margen? (b) ¿Por qué el trigrama con suavizado aditivo puede quedar por detrás del bigrama con estos datos? Relaciónelo con la proporción de casillas vacías y con el sobreajuste. (c) ¿Qué perplejidad tendría un modelo perfecto? (d) Convierta la mejor pérdida a bits por carácter y relaciónela con la entropía de la fuente (M06).
  \item \emph{Mejora} (necesaria para la calificación de excelente). Implemente la interpolación $P=\lambda P_3+(1-\lambda)P_2$, con $P_3$ el trigrama sin suavizar y $P_2$ el bigrama suavizado, y elija $\lambda$ en validación (M01, §3.4). Opcionalmente, añada el modelo de unigramas como referencia intermedia.
\end{enumerate}

\paragraph{Valores de referencia.} Con la partición indicada, y solo para comprobar el código: uniforme, $28$; unigramas, $\approx 15{,}7$; bigramas con $\alpha=0{,}2$, $\approx 12{,}4$; trigramas con $\alpha=0{,}2$, $\approx 13{,}2$; interpolación con bigrama de $\alpha=0{,}2$ y $\lambda=0{,}3$, $\approx 11{,}7$. El tensor de trigramas solo tiene 840 casillas no nulas de las $21\,952$.

\section*{Criterios de evaluación}
\begin{center}
\small
\renewcommand{\arraystretch}{1.15}
\begin{tabularx}{\textwidth}{@{}>{\raggedright\arraybackslash}p{2.3cm}c>{\raggedright\arraybackslash}X>{\raggedright\arraybackslash}X>{\raggedright\arraybackslash}X@{}}
\toprule
Criterio & Peso & Excelente & Suficiente & Insuficiente \\
\midrule
Corrección de los modelos & 30\,\% & Los tres modelos reproducen los valores de referencia; hay \texttt{assert} de normalización y del número de predicciones. & Valores correctos, sin comprobaciones explícitas. & Algún modelo mal normalizado, pérdidas infinitas o relleno incoherente. \\
\addlinespace
Metodología de evaluación & 20\,\% & $\alpha$ elegido en validación y medida final en una partición de prueba independiente. & $\alpha$ elegido en validación. & Se evalúa en entrenamiento o se elige $\alpha$ a la vista de la medida final. \\
\addlinespace
Resultados y visualización & 15\,\% & Tabla completa en nats, bits y perplejidad; gráfica con ejes rotulados y escala logarítmica. & Tabla y gráfica legibles, con alguna carencia. & Falta la tabla o la gráfica. \\
\addlinespace
Análisis & 25\,\% & Responde a (a)--(d) con cifras, explica la dispersión del tensor de trigramas e implementa la interpolación. & Responde a (a)--(d) con cifras. & Respuestas vagas o sin números. \\
\addlinespace
Código y reproducibilidad & 10\,\% & Funciones reutilizables, semillas fijas, ejecución en menos de un minuto. & Se ejecuta de principio a fin. & Errores o dependencia del orden manual de ejecución. \\
\bottomrule
\end{tabularx}
\end{center}
